% Reconstructed text from V5 PDF page 25; non-native source
\par\smallskip\noindent\rule{\textwidth}{0.4pt}\par{\sffamily\footnotesize Original V5 page 25 -- recovered text}\par\begin{flushleft}\footnotesize\ttfamily
\par\smallskip
\par\smallskip
\noindent Diagnostics and Serving Infrastructure for MAS 25\par
\par\smallskip
\par\smallskip
\par\smallskip
\noindent dynamic control flow (conditional branches, native function calls) must still execute client-side\par
\noindent (§6, OSDI 2024).\par
\noindent  DistServe [50] disaggregates prefill and decoding into separate GPU pools connected via\par
\noindent pull-based RDMA, allowing each phase to be independently parallelized and scaled. For a 66B-\par
\noindent parameter model with 512-token requests, each request’s KV cache is approximately 1.13 GB\par
\noindent (§3.3), making transfer cost a first-class design concern. The key architectural insight is that pre-\par
\noindent fill (compute-bound) and decoding (memory-bandwidth-bound) have complementary resource\par
\noindent profiles, so co-locating them on the same GPU pool wastes resources in both directions. By\par
\noindent disaggregating, DistServe achieves 2× goodput improvement on Llama-70B under SLA con-\par
\noindent straints (§5). However, the KV transfer mechanism is intra-request (prefill-to-decode handoff\par
\noindent within a single user’s generation), not inter-agent sharing. No eviction policy is described; KV\par
\noindent state is consumed and released per-request (OSDI 2024). For multi-agent scenarios, DistServe’s\par
\noindent disaggregation could be composed with inter-request prefix sharing (as SGLang or Mooncake\par
\noindent provide), but no such composition has been evaluated.\par
\par\smallskip
\noindent 5.2 Exact KV-Cache Reuse\par
\noindent Exact KV-cache reuse retains computed attention tensors across requests that share identical\par
\noindent token prefixes, avoiding redundant prefill computation. This is the dominant category among\par
\noindent surveyed systems.\par
\par\smallskip
\noindent Foundation: PagedAttention. vLLM [14] introduced PagedAttention, managing KV cache\par
\noindent in non-contiguous fixed-size blocks with a page table. It uses all-or-nothing eviction when\par
\noindent GPU memory is exhausted, preempting entire sequences rather than individual blocks. This\par
\noindent establishes the baseline architecture that most subsequent systems extend.\par
\par\smallskip
\noindent Prefix-tree indexing. SGLang [49] organizes cached KV prefixes in a radix tree with LRU\par
\noindent leaf-first eviction and reference counting. Each node maintains a reference counter; a node is\par
\noindent evictable only when its counter reaches zero (§3). Evicting leaves first enables reuse of common\par
\noindent ancestors until those ancestors themselves become leaves. A key design decision is the dynamic\par
\noindent shared pool: there is no fixed cache budget — “we do not preallocate a fixed-size memory\par
\noindent pool as a cache; instead, we let the cached tokens and the currently running requests share the\par
\noindent same memory pool” (§3). When enough waiting requests exist, the system evicts all cached\par
\noindent tokens in favor of a larger batch size, prioritizing throughput over cache hit rate. Cache-aware\par
\noindent scheduling sorts requests by matched prefix length; Theorem 3.1 proves this achieves optimal\par
\noindent cache hit rate equivalent to DFS order on the radix tree. On Chatbot Arena workloads, SGLang\par
\noindent achieves 52.4\% cache hit rate on LLaVA-NeXT-34B and 74.1\% on Vicuna-33B, with 1.7×\par
\noindent TTFT reduction (§5.5). Across benchmarks, it delivers up to 6.4× throughput improvement\par
\noindent and 3.7× latency reduction over baselines including Guidance v0.1.8 and vLLM v0.2.5 (§5).\par
\noindent The overhead of RadixAttention itself is under 0.3\% of total inference time (§5).\par
\noindent  For our cross-layer analysis, two properties of SGLang are noteworthy. First, the radix\par
\noindent tree stores only KV tensors — no adapter weights, LoRA parameters, or diagnostic state are\par
\noindent managed (§3). The system is single-model per instance, with no multi-adapter or multi-tenant\par
\noindent cache partitioning. Second, the “evict all cached tokens” pressure response is the most aggres-\par
\noindent sive eviction policy among surveyed systems: when memory pressure is high enough, the entire\par
\noindent diagnostic-irrelevant cache is cleared. This design is optimal for request-level throughput but\par
\noindent would need modification to protect diagnostically valuable prefixes. Stated limitations include\par
\end{flushleft}
